\chapter*{\large ВВЕДЕНИЕ}
\addcontentsline{toc}{chapter}{ВВЕДЕНИЕ}

Игра <<Змейка>> хорошо подходит для изучения обучения с подкреплением. В ней каждое действие сразу влияет на дальнейшую ситуацию на поле и на итоговый счёт. Поэтому на такой задаче удобно разбирать, как задаются состояния, действия и награды, и как агент постепенно учится принимать более удачные решения.

Тема работы актуальна потому, что обучение с подкреплением позволяет строить агентов, которые учатся на собственном опыте, а не по заранее прописанным правилам для каждой ситуации. Даже на сравнительно простой игре можно проследить все основные этапы разработки такой системы: создание среды, выбор признаков состояния, настройку награды, обучение и проверку результатов.

Цель курсовой работы --- разработать игру <<Змейка>> на языке Python, реализовать алгоритм обучения с подкреплением для управления агентом и оценить, насколько хорошо он обучается.

Для достижения этой цели были поставлены следующие задачи:
\begin{itemize}
    \item изучить основные идеи обучения с подкреплением и алгоритма Q-learning;
    \item разобрать требования к игровой среде и реализовать логику игры <<Змейка>>;
    \item задать пространство состояний, множество действий и функцию награды;
    \item реализовать агента и процесс накопления Q-таблицы;
    \item провести вычислительные эксперименты и оценить, как меняется качество игры во время обучения;
    \item подготовить таблицы, графики и изображения, которые объясняют устройство программы и поведение агента.
\end{itemize}

Объект исследования --- программная игра <<Змейка>>. Предмет исследования --- методы обучения с подкреплением для выбора действий в дискретной игровой среде. В работе используются методы объектно-ориентированного программирования, имитационного моделирования и вычислительного эксперимента.
